\documentclass[UTF8,11pt,a4paper]{ctexart}
\usepackage[margin=23mm]{geometry}
\usepackage{amsmath,amssymb,booktabs,longtable,array}
\usepackage[colorlinks=true,urlcolor=blue,linkcolor=blue,citecolor=blue]{hyperref}
\usepackage{xurl}
\setlength{\emergencystretch}{3em}
\setlength{\tabcolsep}{4pt}
\title{18号赛题：COPY-only可行性桥梁、闭合证书与矩阵前沿}
\author{R110--R122独立研究记录}
\date{2026年10月11日；旧v14、R102与R109报告分别冻结}
\begin{document}
\maketitle
\begin{abstract}
本轮研究原17个未产出结果的输入，保持原节点各一次、原有序操作数及架构寄存器副作用。
将物理距离窗口可行性化为原顺序的completed/live状态：满窗发射必须最后消费至少一个旧身份，合法顺序可以通过有限COPY旋转物化。
独立原JSON检查的9份闭合证书、6份rook minor证书将15例判为声明身份模型内不可行；另两例未知。
原物理年龄oracle的2000对照和新240阴性对照无矛盾；可选CP-SAT区间模型再对照2240次，构造2118份可重放输出、122个backend否定。
两个矩阵实例各单worker30秒仍UNKNOWN。没有新增原336合法输出，本轮成果是可行性辨别与证明，不是COPY质量改善或官方Level通过。
\end{abstract}

\section{规范、已有质量和本轮问题}
\href{https://www.chaspark.com/\#/races/competitions/1326690549003280384}{官网18题正文}
仍只授权原DAG重排及COPY插入，距离窗口$W=2^b$，$b=2,3,4,5$，引用$t\#W$合法。
正文列Level 0、1、3，Level 2、COPY分母及平均权重未定义，没有额外GET、重算、spill、等值替换或删节点授权。
本文采用原GET/SET/ADD/MULW各一次、每个原GET身份不同、算术原有序参数不变的模型；同寄存器RAW/WAR/WAW限制顺序。
两次GET数值相同不自动授权将原身份合并，带重复寄存器的否定证据尤其需要这一作用域。

R109在原336合成图的运行记录为272合法、47已有身份模型证书、9预算失败、8启发失败；合法输出共192929 COPY。
共享归约年龄缓存带来的质量收益另见冻结的\texttt{project18-age-research.tex}。
本轮不重写这些运行计数：新增15份证据是对原17未输出输入的\emph{离线诊断}，不是默认调度器已经在线生成的新状态，也没有新合法输出。
全部336输入仍保留，不能把成功子集均值替代全分布的评分。
此前必要COPY下界在全部336等权输入为11.7792\%，在既有272合法输入为10.0864\%，均限严格模型及本地分布。
它们不说明未知官方输入不能达到Level 1，也不把不可行图假定为官方有效block。

\section{从物理年龄到原顺序的等价桥梁}
Sugita等的距离ISA代码生成\cite{sugita}说明move与生存区间之间的关系，但其完整代码生成允许spill，不能借用完备性声称本题皆可行。
本轮单独证明严格原身份COPY-only模型的桥梁。
固定一个原指令前缀，$L$是已产生且仍被未来\emph{原数据操作数}需要的不同身份集合。
RAW/WAR/WAW排序边不使临时值额外存活。
一个原排列可物化为合法输出，当且仅当每次原发射满足：
\begin{enumerate}
\item 发射后未来必需身份数量$|L'|\leq W$；
\item 发射前$|L|=W$时，本条原指令至少最后消费一个现存身份。
\end{enumerate}
第二项不能被第一项替代：死GET或SET没有未来输出，却同样推进一个物理槽；满窗且没有最后消费时，会丢掉一个未来必需身份。

必要性来自任意合法输出在原前缀之间必须保留所有$L$于最后$W$个物理槽。
满窗时每槽恰为一个不同身份，COPY只能把最老身份移到最新槽；原发射若不最后消费旧身份，必永久丢失一个后续所需身份。
严格模型不允许重算、重新GET、别名或删除来恢复它。
充分性通过旋转构造：若最老槽的身份已死或被当前原指令最后消费，直接原发射；否则COPY最老必需身份并继续。
满窗最多$W-1$次旋转找到当前last-use；未满窗总有死槽，最多$|L|\leq W-1$次旋转找到可覆盖槽。
每个原指令因此至多新增$W-1$ COPY；这是存在性上界，不是最少COPY算法。

\section{有限状态搜索和可重放否定证据}
\texttt{retention\_search.py}以已完成原节点mask为精确可行性状态。
它唯一确定未来数据生存集合及依赖ready；物理年龄差异可由上述旋转消除，因此不需保留全部COPY历史。
使用迭代有界DFS避免4096节点递归栈，状态、保存选择及时间皆有上限。
合法ready且没有未来数据用途的原节点可优先执行：它只缩小前沿或解锁依赖，不增加未来数据生存，因此该可行性支配不丢可行解。
这不保持最少COPY，不能用于COPY最优证书。
原COPY和不支持opcode拒绝此匹配器；限额耗尽为UNKNOWN，只有完整状态闭合才证明不可行。

闭合证书保存所有搜索mask、真实原依赖和模型参数。
独立\texttt{retention\_certificate.py}重新解析原JSON、有序引用及寄存器hazard，检查初始mask存在、完整mask不存在、每个状态满足依赖闭合和容量、每个合法后继闭合。
若存在支配的死ready节点，则至少保留一个对应后继；不能自行添加原JSON没有的固定顺序约束。
去除初始状态、删后继或篡改hazard依赖均拒绝。
根代理另外未导入项目模块重建9份证书，最大12373个mask，9份通过且9次去闭合篡改拒绝。

原17失败中8个稀疏共享/生存图与1个跨组共享图在此搜索得到闭合证书；另8矩阵实例耗尽1000ms仍未知。
这9份证据没有产出新输出，也不能推广到允许身份合并的语义。

\section{矩阵的原数据图minor下界}
矩阵归约中，原A行GET和B列GET的消费者足迹组成$d\times d$输出格点。
严格识别器从足迹而非family、寄存器名字或种子识别形态，要求所有共享算术拒绝、每格私有算术连通、每行/列足迹有足够不同原GET。
每格分支集合使用该格私有MUL/ADD/SET，以及独占分配的一个行GET与一个列GET。
分支集合必须不相交、内部连通，所有同一行/列相邻格间必须有\emph{真实原数据边}；额外边允许，缺边拒绝。
不能把寄存器hazard排序边加入minor来证明临时窗口下界。

所得minor是rook graph $L(K_{d,d})$。
Harvey--Wood的Theorem 3取$k=2,c=d$给出\cite{harvey}
\[
\operatorname{tw}(L(K_{d,d}))=\frac{d(d+1)}2-1.
\]
合法严格调度的原数据图有前沿路径分解，袋大小至多$W+1$，故其treewidth至多$W$。
经raw分支集合验证的minor下界大于$W$足以排除可行输出；下界小于等于$W$只是不排除，绝不证明可行。

\begin{center}
\begin{tabular}{rrrrl}
\toprule
实际原节点 & $d$ & 必要界 & 测试$W$ & 结论\\
\midrule
160 & 4 & 9 & 8 / 16 & 8不可行，16未知\\
784 & 7 & 27 & 8 / 16 / 32 & 8、16不可行，32未知\\
3744 & 12 & 77 & 8 / 16 / 32 & 全部不可行\\
\bottomrule
\end{tabular}
\end{center}
R114共17输入，8个实际minor中6个得到不可行证据、2个仅必要界；其余9个不匹配而无minor证书。
根代理独立原数据边检查8份minor及8次重叠篡改拒绝。
连同闭合证书，原17现为15份作用域明确的否定证据、2个未知。

\section{可选CP-SAT模型与独立差分}
Chen的min-reg CP-SAT工作\cite{chen}启发了原位置与生存区间建模，但本模型额外处理物理满窗死槽条件。
设原指令位置$p_v$是$0,\ldots,N-1$的全排列，所有原数据及内存依赖满足$p_u<p_v$。
对有后续数据用途的身份$v$，设$q_v=\max_{u\in\mathrm{dataUsers}(v)}p_u$，区间$[p_v,q_v)$以单位需求参与容量$W$的Cumulative。
整数时间$t$表示原指令位置$t$发射\emph{之后}：结果出生即计数，最后消费后源身份不再计数。
无后续数据用途的原节点另设单位区间$[p_v,p_v+1)$，表示仍须容纳死物理槽。
于是每个整数时点恰为after-live加必要死槽，等价于桥梁的两个条件。

模型允许所有原MUL/ADD内部交错，没有把整个输出格点当不可拆原子，也没有generator-specific对称约束。
先祖/后代数仅缩小合法position域；重复操作数去重计算消费者但发射仍恢复有序真实原args。
求解器只求合法原顺序，没有COPY目标。
\texttt{OPTIMAL}表示纯可行性模型求解完成，不表示最少COPY；\texttt{INFEASIBLE}是backend在声明模型的结论，没有便携独立闭合证书；\texttt{UNKNOWN}保持未知。
OR-Tools安装于独立环境，锁定版本9.15.6755及全部11个直接/传递依赖，默认标准库环境和auto不改变\cite{ortools}。

\begin{longtable}{p{.15\textwidth}p{.30\textwidth}p{.29\textwidth}p{.15\textwidth}}
\toprule
实验 & 数量与独立参考 & 实际结果 & 墙时观察\\
\midrule
R111 & 1000输入，自由/固定2000物理年龄Dijkstra & 1962可行、30否定、8参考未知；retention 1970可行、30否定；1992共同完成无矛盾 & 8.565s\\
R112 & 120新hazard屏障/mandatory spine，自由/固定240参考 & 140可行、92否定、8参考未知；retention 148可行、92否定；232共同完成无矛盾 & 9.416s\\
R116 & 同R111全部2000模型，不重新选择输入 & CP 1970构造、30否定、0矛盾 & 46.483s\\
R117 & 同R112全部240模型 & CP 148构造、92否定、0矛盾 & 6.052s\\
R118 & 原160/W16、784/W32各单worker30s & 两个UNKNOWN，0新输出 & 61.644s\\
\bottomrule
\end{longtable}
R112自由顺序120例包含46否定（36同寄存器hazard、10纯mandatory spine），不是仅固定原顺序导致的阴性。
两种物理oracle均穷举合法原发射和活跃COPY动作，有限预算未知保留。
CP的2118份构造都保存实际输出和原顺序，原JSON身份排列、有序参数、符号及8位有符号/16位无符号/32位有符号回绕整数重放通过。
三个整数配置为本地测试语义，不是官方运算位宽。
五个新增边界测试涵盖满窗死GET、hazard死SET、满窗COPY旋转、重复参数/epoch、预算/原COPY拒绝；首次运行另加载五个已有桥梁测试，10项通过。

160节点CP模型30秒搜索131817 branches、45011 conflicts；784节点为24925 branches、840 conflicts。
UNKNOWN不代表原图无解，backend也没有给出新见证。
更大的输入及COPY最优化尚未由此模型完成，不能以小图正确率冒充大图质量。

\section{资源图对称消除及保留的无效实验}
直接观察两个矩阵的共享行/列，进一步构建可选对称约束。
首先要求纯结构、GET和SET寄存器唯一且读写集合分离；随后raw验证每个相邻行交换与列交换是覆盖\emph{全部原节点}的资源数据图置换。
生成元逐节点保留op及数据参数多重集合；行交换只交换对应行GET组、逐列保持列GET组集合，列交换同理。
不通过名字、family或minor存在性推断automorphism，缺边、缺生成元、组重叠或hazard输入均拒绝。

令每行首次GET位置$b_r=\min_{v\in\text{row }r}p_v$。
任何可行顺序在已验证生成元作用下仍保持before/after live和last-consume条件，因此可以按$b_r$排序行；该置换逐列保持列GET集合，不改变列首次出生位置。
再排序列，逐行保持行组，不破坏行排序。
于是加入两个轴的首次出生位置递增约束不丢声明模型的可行性；这里只证明可行性，不保证最少COPY或求解时间。
原MUL/ADD内部顺序不额外固定，$k$层不统一，$C_{ij}$不原子化。
实际发射仍使用真实原ID、reg、op和有序args，资源置换不是等值替换或交换操作数授权。

R120新24图覆盖$d=2,3$、四个窗口、三个种子，重新随机拓扑排列，使用opaque原ID/寄存器及逐算术有序参数反转。
与实际冻结R119原CP模型同worker/seed、每例2s配对：两侧均9构造、3backend否定、12UNKNOWN，0矛盾。
18个输出独立身份/符号/三整数重放通过。
共同9图COPY合计63到61，但两例13到16、1到2退化；这是没有COPY目标的可行性模型，不能视为默认质量改善。
24图48次总墙时51.255s，72节点的12例在2s大多停于预处理、无搜索branches，说明短预算模型成本本身不可忽略。
首次小图测试误预期24节点/W4应可行，原CP否定与独立retention闭合227mask一致；测试改为明确阴性，没有修改算法迎合预期。
新的四项对称边界加五项CP边界共9项回归2.785s通过。

R121对原两个矩阵使用同一30s单worker预算，真正冻结R119源码与新对称模型配对；实际.pbtxt模型与原JSON哈希保存。
两侧各2个UNKNOWN，4次整体墙时121.500s，没有产生新输出；完整分类及代价以\texttt{experiments/r121-cpsat-symmetry-matrix/report.json}保留。
对称约束并未消除这两例困难，不继续仅扩大超时宣称成果。

\section{资源、成本、冻结及后续余量}
本轮实际测试在Windows、Ryzen 9 8945HX（16核32线程）、约31.3GiB RAM、CPython 3.12.14运行。
矩阵CP始终1worker顺序执行；宿主其他应用使空闲内存一度约1.6--2.1GiB，没有关闭其他应用。
一次进程观测矩阵实际worker工作集约107MiB，CIM累计峰字段约164MiB；非连续测量，不能当两个矩阵完整峰内存保证。
固定搜索预算覆盖模型构建和求解，首次backend导入、输出、序列化、独立验证另计；不是端到端硬墙时上限。
以上时间为功能批墙时记录，可能与其他项目功能测试重叠，不能宣称隔离性能收益。
默认测试无需GPU/开发板；本轮可选依赖下载约50MiB及其他小依赖，保留独立环境。
最低双核64位CPU/4GiB RAM/5GiB空闲磁盘、建议8核/16GiB/10GiB是规划建议，未低配实机复验，且搜索内存须额外留余量。

原数据及实际源哈希分别冻结于R113、R115、R119、R122目录的\texttt{source-provenance.json}。
R116/R117逐例保存输入hash、实际\texttt{.pbtxt}模型hash、backend版本/seed/budget、输出和原ID见证；R118保存两个原输入及实际模型，不覆盖未知。
R119冻结时间UTC 2026-10-11 08:26:06.833829；旧v14、R102、R109报告不修改。

剩余研究优先级为两个矩阵的实际可行见证或更强原数据下界，以及cross-group/ladder的一般共享DAG COPY差距。
近窗口脊柱占大量COPY，但已有必要界很紧，不仅根据COPY总量重复投入。
对矩阵不能只提高超时取得幸运结果；需要分析行/列根出生、逐$k$的MUL前沿和私有ADD链，允许内部交错，再形成有解释的候选或下界。
对于可行图，completed/live可行性搜索未优化物理年龄，需在保持旧候选的前提下引入具体年龄成本或顺序局部搜索。
一名开发者的后续估算：矩阵结构/状态证明3--5工作日，一般共享COPY候选3--5日，独立持出与低配/跨平台复验2--3日。
生成器、原JSON独立审计、环境资源测量可并行；核心模型及证明须分阶段合并。
这些是工作量估计，不是已耗人日或官方达标承诺。

\begin{thebibliography}{9}
\bibitem{sugita} Sugita et al. \emph{A Sound and Complete Algorithm for Code Generation in Distance-Based ISA}. CC 2023. \url{https://www.rsg.ci.i.u-tokyo.ac.jp/lab/en/assets/files/cc2023-sugita.pdf}.
\bibitem{chen} Gang Chen. \emph{Optimal and Heuristic Min-Reg Scheduling Algorithms for GPU Programs}. \url{https://arxiv.org/abs/2303.06855v4}.
\bibitem{harvey} Daniel J. Harvey and David R. Wood. \emph{Treewidth of the Line Graph of Complete and Complete Multipartite Graphs}, Theorem 3, arXiv:1210.8205v2. \url{https://arxiv.org/abs/1210.8205}.
\bibitem{ortools} Google OR-Tools. Python installation and CP-SAT API. \url{https://developers.google.com/optimization/install/python}; \url{https://or-tools.github.io/docs/pdoc/ortools/sat/python/cp_model.html}.
\end{thebibliography}
\end{document}
